\documentclass[11pt,a4paper]{article}
\usepackage[T1]{fontenc}
\usepackage[utf8]{inputenc}
\usepackage{lmodern}
\usepackage{amsmath,amssymb,amsthm,mathtools,mathrsfs}
\usepackage{graphicx,float,xcolor,multicol}
\usepackage[shortlabels]{enumitem}
\usepackage[margin=27mm]{geometry}
\usepackage{microtype,needspace,placeins}
\usepackage{tikz,tikz-cd}
\usetikzlibrary{arrows.meta,calc,positioning,patterns,decorations.pathreplacing}
\usepackage[colorlinks=true,linkcolor=teal,urlcolor=teal,citecolor=teal]{hyperref}
\setlength{\parindent}{0pt}
\setlength{\parskip}{.6em}
\setcounter{secnumdepth}{0}
\allowdisplaybreaks
\raggedbottom
\providecommand{\cross}{\times}
\DeclareUnicodeCharacter{2020}{\ensuremath{\dagger}}
\DeclareMathOperator{\supp}{supp}
\DeclareMathOperator*{\esssup}{ess\,sup}
\providecommand{\chapter}{\section}
\newenvironment{tool}[2]{\subsection*{Tool #1 --- #2}}{}
\newcommand{\ToolRef}[1]{Tool #1}
\newcommand{\FigureTag}[2]{\textit{Figure #1}}
\newcommand{\Result}[1]{\subsection*{#1}}
\newcommand{\Application}[1]{\subsubsection*{Application\if\relax\detokenize{#1}\relax\else\space--- #1\fi}}
\newcommand{\ProofHeading}{\par\textit{Proof.}\space}
\newcommand{\ProofOf}[1]{\par\textit{Proof of #1.}\space}
\newcommand{\RemarkHeading}{\par\textbf{Remark.}\space}
\newcommand{\NoteHeading}{\par\textbf{Note.}\space}
\newcommand{\ExerciseHeading}{\par\textbf{Exercise.}\space}

% Rendering support only; mathematical text is in each independent post.
\usepackage{booktabs,array,longtable}
\definecolor{HandoutBlue}{HTML}{2F6690}
\definecolor{HandoutNavy}{HTML}{17365D}
\newtheorem{theorem}{Theorem}
\newtheorem{lemma}[theorem]{Lemma}
\newtheorem{proposition}[theorem]{Proposition}
\newtheorem{corollary}[theorem]{Corollary}
\newtheorem{conjecture}[theorem]{Conjecture}
\newtheorem{definition}{Definition}
\newtheorem{example}{Example}
\newtheorem{namedtoolcounter}{Tool}
\newenvironment{namedtool}[1]{\begin{namedtoolcounter}[#1]}{\end{namedtoolcounter}}
\newtheorem*{claim}{Claim}
\newtheorem*{remark}{Remark}
\newtheorem*{exercise}{Exercise}
\newtheorem*{question}{Question}
\newtheorem*{observation}{Observation}
\newcommand{\SourceChapter}[2]{\section*{#2}}
\newcommand{\Lecture}[2]{\section*{Lecture #1: #2}}
\newcommand{\unclear}[1]{\textit{[unclear: #1]}}
\newcommand{\sourcegap}{\textit{[unfinished in the source]}}
\DeclareMathOperator{\dist}{dist}
\DeclareMathOperator{\id}{id}
\DeclareMathOperator{\Hol}{Hol}
\DeclareMathOperator{\Per}{Per}
\DeclareMathOperator{\Fix}{Fix}
\DeclareMathOperator{\diam}{diam}
\DeclareMathOperator{\Var}{Var}
\DeclareMathOperator{\Leb}{Leb}
\DeclareMathOperator{\Hom}{Hom}
\DeclareMathOperator{\End}{End}
\DeclareMathOperator{\Ann}{Ann}
\DeclareMathOperator{\Spec}{Spec}
\DeclareMathOperator{\Max}{Max}
\DeclareMathOperator{\Ob}{Ob}
\DeclareMathOperator{\Tor}{Tor}
\DeclareMathOperator{\Ass}{Ass}
\DeclareMathOperator{\Mor}{Mor}
\DeclareMathOperator{\Fun}{Fun}
\DeclareMathOperator{\Nat}{Nat}
\DeclareMathOperator{\Cone}{Cone}
\DeclareMathOperator{\MCG}{MCG}
\DeclareMathOperator{\Aut}{Aut}
\DeclareMathOperator{\Deck}{Deck}
\DeclareMathOperator{\SL}{SL}
\DeclareMathOperator{\PSL}{PSL}
\DeclareMathOperator{\Area}{Area}
\DeclareMathOperator{\im}{im}
\DeclareMathOperator{\PSh}{PSh}
\newcommand{\CRing}{\mathsf{CRings}}
\newcommand{\AMod}{A\text{-}\mathsf{Mod}}
\newcommand{\Ab}{\mathsf{Ab}}
\newcommand{\Z}{\mathbb Z}
\newcommand{\Q}{\mathbb Q}
\newcommand{\R}{\mathbb R}
\newcommand{\C}{\mathbb C}
\newcommand{\N}{\mathbb N}
\newcommand{\kk}{\mathbb k}
\renewcommand{\aa}{\mathfrak a}
\newcommand{\bb}{\mathfrak b}
\newcommand{\pp}{\mathfrak p}
\newcommand{\qq}{\mathfrak q}
\newcommand{\mm}{\mathfrak m}
\newcommand{\nn}{\mathfrak n}
\newcommand{\rad}{\mathrm r}
\newcommand{\cat}[1]{\mathsf{#1}}
\setcounter{secnumdepth}{0}
\allowdisplaybreaks

\graphicspath{{figures/lemma-book/}}
\begin{document}
\section*{Problems Proposed by S. Ramanujan}
% BLOG-CONTENT-BEGIN
Here is a collection of problems proposed by S.~Ramanujan.
\begin{enumerate}[label=\arabic*.]
\setcounter{enumi}{0}\item Prove without calculus that
$\frac32\log2=1+\frac2{4^3-4}+\frac2{8^3-8}+\frac2{12^3-12}+\cdots$.
\setcounter{enumi}{1}\item Show that it is possible to solve the equations
\[
\begin{array}{ll}
x+y+z=a,&px+qy+rz=b,\\
p^2x+q^2y+r^2z=c,&p^3x+q^3y+r^3z=d,\\
p^4x+q^4y+r^4z=e,&p^5x+q^5y+r^5z=f.
\end{array}
\]
\setcounter{enumi}{2}\item Find $\sqrt{1+2\sqrt{1+3\sqrt{1+4\sqrt{1+\cdots}}}}$ and
$\sqrt{6+2\sqrt{7+3\sqrt{8+\cdots}}}$.
% Source page 19 - left
\setcounter{enumi}{3}\item If $\alpha\beta=\pi$, show that
$\sqrt\alpha\int_0^\infty e^{-x^2}/\cosh(\alpha x)\,dx
=\sqrt\beta\int_0^\infty e^{-x^2}/\cosh(\beta x)\,dx$.
\setcounter{enumi}{4}\item Show that Euler's constant
$\gamma=\lim_{n\to\infty}(1+1/2+\cdots+1/n-\log n)$ equals
\[
\log2-1\left(\frac2{3^3-3}\right)
-2\left(\frac2{6^3-6}+\frac2{9^3-9}+\frac2{12^3-12}\right)-\cdots,
\]
the first term in the $n$th group being
$2/\bigl(((3^n+3)/2)^3-(3^n+3)/2\bigr)$.
\setcounter{enumi}{5}\item Show that
\[
\prod_{n=1}^\infty\left(1+\frac1{n^3}\right)=\frac1\pi\cosh\frac{\pi\sqrt3}{2},
\qquad
\prod_{n=2}^\infty\left(1-\frac1{n^3}\right)=\frac1{3\pi}\cosh\frac{\pi\sqrt3}{2}.
\]
\setcounter{enumi}{6}\item For any odd positive integer $n$, show that
$\int_0^\infty\sin(nx)/\bigl((\cosh x+\cos x)x\bigr)\,dx=\pi/4$, and hence prove
\[
\frac\pi8=\frac1{\cosh(\pi/(2n))+\cos(\pi/(2n))}
-\frac1{3\bigl(\cosh(3\pi/(2n))+\cos(3\pi/(2n))\bigr)}+\cdots.
\]
% Source page 19 - right
\setcounter{enumi}{7}\item If $\sin(x+y)=2\sin((x-y)/2)$ and $\sin(y+z)=2\sin((y-z)/2)$, prove
\[
\sqrt[4]{\tfrac12\sin x\cos z}+\sqrt[4]{\tfrac12\cos x\sin z}
=\sqrt[12]{\sin2y}.
\]
Verify this when $\sin2x=(\sqrt5-2)^3(4+\sqrt{15})^2$,
$\sin2y=\sqrt5-2$, and $\sin2z=(\sqrt5-2)^3(4-\sqrt{15})^2$.
\setcounter{enumi}{8}\item Show that $\sum_{n=1}^\infty n/(e^{2n\pi}-1)=1/24-1/(8\pi)$.
\setcounter{enumi}{9}\item Show that
\begin{align*}
(6a^2-4ab+4b^2)^3
&=(3a^2+5ab-8b^2)^3+(4a^2-4ab+6b^2)^3\\
&\quad +(5a^2-5ab-3b^2)^3,
\end{align*}
and find other quadratic expressions.
As written, the displayed expression is not a polynomial identity: at
$a=0,b=1$, its two sides are $64$ and $-323$. One of the coefficients in
the source therefore needs correction.
\setcounter{enumi}{10}\item $2^n-7$ is a perfect square for $n=3,4,5,7,15$. Find other values.
\setcounter{enumi}{11}\item Show that
\[
\prod_{j=0}^\infty\bigl(1+e^{-(2j+1)\pi\sqrt{55}}\bigr)
=\left(\frac{1+\sqrt{3+2\sqrt5}}{\sqrt2}\right)
\bigl(e^{-\pi\sqrt{55}}\bigr)^{1/24}.
\]
\end{enumerate}
% BLOG-CONTENT-END
\end{document}
